% ============================================================
%  Section 6 --- Expériences complémentaires
% ============================================================

\section{Expériences complémentaires}
\label{sec:experiments}

Les quatre expériences présentées ici approfondissent les résultats des hypothèses
principales selon quatre axes : la contribution de chaque modalité au clustering (H1)
et à la prédictibilité (H3), la comparaison d'architectures d'encodeur, la structure
causale inter-services révélée par l'ontologie, et la transférabilité du pipeline
sur un dataset externe.

\subsection{Ablation des modalités --- H1 (réentraînement)}
\label{sec:ablation_h1}

\paragraph{Protocole.}
Pour mesurer la contribution de chaque modalité à la structuration de l'espace latent
(H1), chaque condition d'ablation réentraîne complètement l'encodeur et le réseau siamois
sur les features réduites. Cette approche de réentraînement complet est méthodologiquement
rigoureuse : un masquage à l'inférence (OOD) produirait des embeddings hors distribution
et biaiserait systématiquement les scores vers le bas. Sept conditions sont évaluées :
les 3 modalités en isolation (M, T, L) et les 4 paires et triplet possibles.

\paragraph{Résultats.}

\begin{table}[h]
\centering
\caption{Ablation modalités sur H1 --- réentraînement complet (graine 42).}
\label{tab:ablation_h1}
\begin{tabular}{lcccr}
\toprule
\textbf{Condition} & \textbf{$n_\text{feat}$} & \textbf{sil\_train} & \textbf{sil\_test} & \textbf{$\Delta$ vs full} \\
\midrule
\textbf{full}    & 17 & 0.378 & 0.439 & --- \\
\textbf{M\_only} & 7  & 0.241 & \textbf{0.497} & \textbf{+0.058} \\
T\_only          & 6  & 0.064 & 0.412 & -0.027 \\
M+L              & 11 & 0.251 & 0.382 & -0.057 \\
T+L              & 10 & 0.022 & 0.341 & -0.098 \\
M+T              & 13 & 0.318 & 0.316 & -0.123 \\
L\_only          & 4  & -0.138 & 0.051 & -0.388 \\
\bottomrule
\end{tabular}
\end{table}

\paragraph{Interprétation.}
Le résultat le plus contre-intuitif est que M\_only ($n_\text{feat} = 7$) surpasse
le modèle full ($+0.058$ sil\_test). La combinaison M+T est pire que M\_only ($-0.123$),
ce qui indique que les features de traces $T(t)$ \emph{dégradent} la géométrie siamoise
lorsqu'elles sont ajoutées aux métriques sur un jeu d'entraînement de seulement 209 épisodes.

L'explication probable est un sur-paramétrage : le STGCN dispose de suffisamment de
capacité pour ajuster les 17 features, mais les features de traces et de logs introduisent
de la variabilité intra-cluster (un même type de panne peut avoir des signatures de logs
très différentes d'un épisode à l'autre) qui brouille la structure contrastive.
L\_only présente même un sil\_train négatif --- la modalité logs seule ne structure pas
l'espace latent sur les données disponibles.

Ce résultat contraste fortement avec celui de H3 (Section~\ref{sec:ablation_h3}) :
les modalités $T$ et $L$ sont géométriquement bruitées pour le clustering, mais
indispensables pour la prédictibilité des précurseurs.

\subsection{Ablation features --- H3 (masquage à l'inférence)}
\label{sec:ablation_h3}

\paragraph{Protocole.}
Pour H3, la question est différente : quelles features le classifieur LR entraîné
exploite-t-il le plus pour prédire les précurseurs ? Cette mesure de sensibilité
s'effectue par masquage à l'inférence --- les features sont mises à zéro dans les
fenêtres pré-injection $[t-k^*, t]$ sans réentraînement, et l'impact sur l'AUROC macro
est mesuré. Il s'agit explicitement d'une mesure de sensibilité du modèle entraîné,
non d'une importance causale.
Vingt-quatre conditions sont évaluées : 7 ablations par modalité et 17 ablations
leave-one-out (une feature à la fois).

\paragraph{Résultats par modalité.}

\begin{table}[h]
\centering
\caption{Ablation modalités sur H3 --- masquage à l'inférence (graine 42, AUROC macro moyen).}
\label{tab:ablation_h3_modality}
\begin{tabular}{lcc}
\toprule
\textbf{Condition} & \textbf{AUROC macro} & \textbf{$\Delta$ vs full} \\
\midrule
\textbf{full}    & \textbf{0.954} & --- \\
M+L              & 0.916          & -0.038 \\
M+T              & 0.891          & -0.063 \\
M\_only          & 0.756          & -0.198 \\
T+L              & 0.563          & -0.391 \\
T\_only          & 0.520          & -0.434 \\
L\_only          & 0.488          & -0.466 \\
\bottomrule
\end{tabular}
\end{table}

\paragraph{Features les plus critiques (leave-one-out).}

\begin{table}[h]
\centering
\caption{Top features LOO pour H3 --- masquage à l'inférence (graine 42).}
\label{tab:ablation_h3_loo}
\begin{tabular}{lcc}
\toprule
\textbf{Feature} & \textbf{AUROC sans} & \textbf{$\Delta$} \\
\midrule
\texttt{disk\_io}          & 0.866 & -0.088 \\
\texttt{lexical\_entropy}  & 0.905 & -0.049 \\
\texttt{latency\_p99}      & 0.912 & -0.042 \\
\bottomrule
\end{tabular}
\end{table}

\paragraph{Interprétation.}
Le modèle full (AUROC $= 0.954$) surpasse toutes les réductions, ce qui est l'inverse
de l'ablation H1. Cette inversion révèle que les modalités jouent des rôles distincts
selon la tâche : $M(t)$ fournit une géométrie propre pour le clustering contrastif,
tandis que $T(t)$ et $L(t)$ portent des signatures temporelles pré-injection qui ne
contribuent pas à la structure des clusters mais permettent de prédire leur apparition.

Le résultat le plus frappant est la criticité de \texttt{disk\_io} ($\Delta = -0.088$),
malgré ses 16.7\,\% de NaN. Cette feature capture l'augmentation des IOPS précédant
certains types de pannes (saturation de ressource, fuite mémoire), un signal précurseur
d'autant plus informatif qu'il est peu corrélé aux autres métriques.
Son importance réelle est probablement sous-estimée sur ewat\_v3 à cause des NaN
--- argument direct pour résoudre L2 dans ewat\_v4.

Deux paires de features sont fortement redondantes :
\texttt{latency\_p99} $\leftrightarrow$ \texttt{span\_dur\_median} ($\rho = 0.936$) et
\texttt{error\_rate\_http} $\leftrightarrow$ \texttt{abnormal\_span\_rate} ($\rho = 0.927$).
À l'inverse, \texttt{retry\_rate}, \texttt{log\_warn\_rate} et \texttt{queue\_depth}
présentent un $\Delta \approx 0$ en LOO --- ces features sont candidates à la suppression
dans ewat\_v4, ce qui réduirait $S(t)$ de $\mathbb{R}^{N \times 17}$
à $\mathbb{R}^{N \times 14}$.

\subsection{Comparaison d'architectures d'encodeur}
\label{sec:architectures}

\paragraph{Protocole.}
Trois architectures d'encodeur sont comparées sur ewat\_v3 (graine 42), en utilisant
la même méthodologie corrigée (nearest centroid pour H1, $k^*$ sur val pour H3) :
\begin{itemize}
  \item \textbf{STGCN} (baseline) : convolution graphique GCN + TCN temporel.
  \item \textbf{SimCLR} (NT-Xent) : pré-entraînement contrastif auto-supervisé
    \cite{chen2020simclr}, sans labels de scénarios pendant l'encodage.
  \item \textbf{GAT} (Graph Attention Network) : mécanisme d'attention sur les arêtes
    de $G(t)$ \cite{velivckovic2018}, remplaçant la convolution GCN fixe par une
    pondération apprise des voisins.
\end{itemize}

\paragraph{Résultats.}

\begin{table}[h]
\centering
\caption{Comparaison architectures encodeur (graine 42, ewat\_v3).}
\label{tab:architectures}
\begin{tabular}{lcccccc}
\toprule
\textbf{Architecture} & \textbf{$K$} & \textbf{sil\_val} & \textbf{sil\_test} & \textbf{H1} & \textbf{Types H3} & \textbf{AUROC moyen} \\
\midrule
STGCN (baseline) & 10 & 0.470 & 0.414 & \checkmark & 8/10 & 0.954 \\
SimCLR (NT-Xent) & 15 & 0.495 & 0.429 & \checkmark & 11/15 & \textbf{0.964} \\
GAT (attention)  & 15 & 0.445 & \textbf{0.497} & \checkmark & \textbf{13/15} & 0.929 \\
\bottomrule
\end{tabular}
\end{table}

\paragraph{Interprétation.}
Les trois architectures valident H1, ce qui confirme que la structurabilité des types
de pannes est une propriété robuste du signal et non un artefact de l'architecture STGCN.

Le GAT produit la meilleure géométrie de l'espace latent (sil\_test $= 0.497$, $+0.083$
vs STGCN) et le plus grand nombre de types prédictibles ($13/15$). Le mécanisme
d'attention sur les arêtes de $G(t)$ permet d'apprendre dynamiquement quelles relations
inter-services sont informatives pour chaque type de panne, ce qui améliore la
structuration sans nécessiter une topologie fixe.

SimCLR présente le meilleur AUROC moyen ($0.964$) mais 4 types NaN (clusters trop
petits en test, $n_\text{pos} < 2$) --- conséquence d'un $K = 15$ plus élevé qui fragmente
les scénarios rares. Le pré-entraînement NT-Xent améliore la séparabilité globale mais
produit des clusters de taille inégale sur un dataset de 299 épisodes.

STGCN est retenu comme architecture principale du pipeline car $K = 10$ est le plus
stable (écart-type $= 0.7$ sur 5 graines), ses résultats multi-graines sont disponibles,
et le compromis silhouette/AUROC est satisfaisant. Le GAT est le candidat naturel pour
ewat\_v4, où des épisodes plus longs devraient amplifier l'avantage de l'attention
adaptative.

Pour contextualiser la valeur ajoutée du STGCN, les baselines B3 et B4 mesurent
l'AUROC macro sur les scénarios Chaos Mesh (vérité terrain indépendante, $k = 6$) :
B3 (features brutes, sans STGCN) $= 0.835$, B4 (STGCN) $= 0.835$.
Cette coïncidence numérique masque une redistribution non triviale : l'encodeur améliore
\texttt{fail\_slow\_cpu} de $+0.270$ et dégrade \texttt{noisy\_neighbor} de $-0.246$.
La valeur du STGCN est donc géométrique (structuration pour le clustering siamois,
H1 sil $= 0.519$) et redistributive, plutôt que globalement prédictive.

\subsection{Ontologie --- relations causales service-level}
\label{sec:ontology_results}

\paragraph{Protocole.}
L'ontologie est construite en deux passes. La passe cluster-level agrège toutes les
séries temporelles d'un cluster avant d'estimer la Transfer Entropy --- ce qui introduit
un biais écologique (les effets intra-épisode sont confondus avec les effets inter-épisodes).
La passe service-level estime la TE intra-épisode (estimateur KSG hiérarchique),
puis agrège les p-values et les valeurs de TE par cluster. Seule la passe service-level
est reportée comme résultat principal.

\paragraph{Résultats --- cluster-level.}
La passe cluster-level produit 22 relations temporelles et 0 relation causale avec
100 permutations (les 2 relations observées lors d'un dry-run à 20 permutations étaient
des faux positifs, confirmant la nécessité d'un nombre suffisant de permutations).

\paragraph{Résultats --- service-level.}

\begin{table}[h]
\centering
\caption{Relations causales service-level par cluster (Transfer Entropy KSG, $p < 0.05$, BH-FDR, 100 permutations).}
\label{tab:ontology}
\begin{tabular}{llcc}
\toprule
\textbf{Cluster} & \textbf{Scénario dominant} & \textbf{Relations} & \textbf{Relation caractéristique} \\
\midrule
C0 & fail\_slow\_cpu        & 23 & --- \\
C1 & drift\_traffic\_ramp   & 19 & --- \\
C2 & resource\_leak         & 8  & $\text{load-gen} \to \text{frontend}$ (TE=0.187) \\
C3 & noisy\_neighbor        & 10 & --- \\
C4 & crash                  & 20 & --- \\
C5 & drift\_rolling\_deploy & \textbf{0}  & --- \\
C6 & drift\_config\_change  & \textbf{0}  & --- \\
C7 & cpu\_starvation        & 15 & --- \\
C8 & faulty\_deploy\_overlap & 20 & $\text{cart} \to \text{load-gen}$ (TE=0.031) \\
C9 & drift\_scale\_up        & 9  & --- \\
\midrule
\textbf{Total} & & \textbf{124} & \\
\bottomrule
\end{tabular}
\end{table}

\paragraph{Interprétation.}
Le résultat le plus significatif est l'absence totale de relations causales pour C5
(\texttt{drift\_rolling\_deploy}) et C6 (\texttt{drift\_config\_change}) : les drifts
bénins ne produisent pas de cascade causale mesurable entre services.
Ce résultat est scientifiquement attendu --- un déploiement progressif modifie la
distribution du signal sans créer de couplage anormal entre services --- et il valide
rétrospectivement la partition drift/anomalie de la formalisation.

La relation $\text{load-generator} \to \text{frontend}$ est ubiquitaire (présente dans
8/10 clusters actifs), reflet du couplage structurel de trafic : le load-generator
drive directement la charge frontend.
Son amplitude varie de TE $= 0.047$ (C4, crash --- effet masqué par l'interruption
brutale du service) à TE $= 0.187$ (C2, resource\_leak --- l'épuisement progressif
des ressources amplifie les délais de propagation entre services).

La relation $\text{cart} \to \text{load-generator}$ est unique à C8
($\theta_{\text{drift} \cap \text{anomaly}}$). Ce sens de causalité --- le service
cart influençant le load-generator --- est cohérent avec les dynamiques de retry/backoff
pendant un déploiement défectueux : les erreurs du cart déclenchent des retries
qui augmentent la charge apparente mesurée par le load-generator.

\subsection{Transfert --- RCAEval RE2-OB}
\label{sec:rcaeval}

\paragraph{Contexte.}
RCAEval RE2-OB est un dataset public composé de 90 épisodes ($30$ types de pannes
$\times 3$ instances) collectés sur \emph{Online Boutique} déployé sur un cluster
Kubernetes distinct. Il partage les 6 services EWAT et les mêmes types de métriques,
mais diffère sur deux points critiques : les épisodes durent 48 steps (vs 21 pour ewat\_v3,
soit un facteur $2.3$) et les scénarios sont différents des injections Chaos Mesh d'ewat\_v3.

\subsubsection{Zero-shot}
\label{sec:rcaeval_zeroshot}

Quatre stratégies de normalisation sont testées sans réentraînement, en appliquant
directement le pipeline ewat\_v3 (encodeur, scaler, centroides) sur les épisodes RCAEval.

\begin{table}[h]
\centering
\caption{Transfert zero-shot --- RCAEval RE2-OB (4 stratégies de normalisation).}
\label{tab:rcaeval_zeroshot}
\begin{tabular}{llcc}
\toprule
\textbf{Stratégie} & \textbf{Features} & \textbf{sil H1} & \textbf{AUROC H3} \\
\midrule
ewat\_v3 scaler  & 17 & 0.778\,* & 0.510 \\
rcaeval scaler   & 17 & 0.234    & 0.497 \\
instance norm    & 17 & 0.287    & 0.507 \\
\textbf{instance norm} & \textbf{M(t) seul} & \textbf{0.684} & \textbf{0.495} \\
\bottomrule
\end{tabular}
\end{table}

La stratégie ewat\_v3 scaler produit un sil $= 0.778$ artificiellement élevé : la
normalisation ewat\_v3 concentre tous les épisodes RCAEval dans une région restreinte de
l'espace latent, ce qui donne une silhouette élevée pour un seul cluster dominant
(81/90 épisodes mappés sur C2, \texttt{resource\_leak}) plutôt que pour une vraie structure.

La meilleure stratégie (instance norm + M(t) seul) atteint H1 PASS (sil $= 0.684$)
mais H3 FAIL (AUROC $\approx 0.5$) : l'encodeur détecte qu'une anomalie est présente,
mais ne peut pas discriminer les types car les classifieurs precurseurs entraînés sur
les fenêtres de 21 steps ewat\_v3 ne reconnaissent pas les signatures à horizon variable
sur des épisodes de 48 steps.

\subsubsection{Few-shot --- Stratégie A}
\label{sec:rcaeval_fewshot}

\paragraph{Protocole.}
La Stratégie A réajuste uniquement le StandardScaler sur $n_\text{few}$ épisodes
RCAEval labellisés, en conservant l'encodeur et les classifieurs ewat\_v3 inchangés.
La boucle couvre $n_\text{few} \in \{1, 3, 5, 10, 20, 40\}$ avec 5 répétitions
aléatoires (IC 95\,\% bootstrap). Le hold-out test correspond aux épisodes restants
($\approx 80$\,\%), stratifié par type de panne.

\begin{table}[h]
\centering
\caption{Transfert few-shot --- Stratégie A (refit scaler, $n_\text{repeats}=5$).}
\label{tab:rcaeval_fewshot}
\begin{tabular}{ccc}
\toprule
\textbf{$n_\text{few}$} & \textbf{sil H1} & \textbf{AUROC H3} \\
\midrule
1  & $0.29 \pm 0.05$ & $0.49 \pm 0.03$ \\
3  & $0.34 \pm 0.04$ & $0.50 \pm 0.02$ \\
5  & $0.38 \pm 0.03$ & $0.50 \pm 0.02$ \\
10 & $0.44 \pm 0.02$ & $0.50 \pm 0.01$ \\
20 & $0.42 \pm 0.02$ & $0.50 \pm 0.01$ \\
40 & $0.40 \pm 0.01$ & $0.50 \pm 0.01$ \\
\bottomrule
\end{tabular}
\end{table}

\paragraph{Interprétation.}
La Stratégie A améliore H1 pour $n_\text{few} \leq 10$ (sil $\approx 0.3$--$0.44$,
PASS au seuil $0.3$), mais H3 reste bloqué à AUROC $\approx 0.50$ quel que soit
$n_\text{few}$.

L'adaptation du scaler corrige partiellement le domain shift de distribution
(d'où l'amélioration de H1), mais les classifieurs LR ewat\_v3 ne peuvent pas
discriminer les types de pannes RCAEval pour deux raisons complémentaires :
(1) la durée des épisodes ($\times 2.3$) change la position temporelle des signatures
pré-injection relative à l'horizon $k^*$ appris sur ewat\_v3 ;
(2) les types de pannes RCAEval ne correspondent pas aux 10 clusters ewat\_v3,
donc les classifieurs one-vs-rest ne trouvent pas de signal positif.

La Stratégie B --- fine-tuning des classifieurs LR ou de la dernière couche de
l'encodeur avec quelques épisodes labellisés RCAEval --- est la prochaine étape
nécessaire pour atteindre H3 AUROC $> 0.7$ en transfert.
